% ECONOMICS_PROBLEM: {"id": "C73-63", "title": "Minimax self-play regret under adversarially valid decentralized learning", "jel_code": "C73", "source_id": "OP-0147"}
% !TeX program = xelatex
\documentclass[11pt,a4paper]{article}
\usepackage[margin=23mm]{geometry}
\usepackage{fontspec}
\setmainfont{Latin Modern Roman}
\usepackage{amsmath,amssymb,mathtools}
\usepackage{xcolor,enumitem,booktabs,longtable,array,xurl,hyperref}
\definecolor{muted}{HTML}{53636C}
\hypersetup{colorlinks=true,urlcolor=blue,linkcolor=blue}
\setlength{\parindent}{0pt}
\setlength{\parskip}{5pt}
\newcommand{\R}{\mathbb R}
\newcommand{\E}{\mathbb E}
\newcommand{\N}{\mathbb N}
\newcommand{\ind}{\mathbf 1}
\newcommand{\Var}{\operatorname{Var}}
\newcommand{\Cov}{\operatorname{Cov}}
\newcommand{\supp}{\operatorname{supp}}
\DeclareMathOperator*{\argmin}{arg\,min}
\DeclareMathOperator*{\argmax}{arg\,max}
\newcommand{\entrymeta}[3]{{\sffamily\small\color{muted}\textbf{\texttt{#1}}\quad|\quad #2\par #3\par}}
\newcommand{\Problemref}[1]{\texttt{#1}}
\newcommand{\JELref}[2]{\texttt{#2} (JEL \texttt{#1})}
\begin{document}
\section*{Minimax self-play regret under adversarially valid decentralized learning}
\textbf{Problem ID:} \texttt{C73-63}\quad \textbf{JEL:} \texttt{C73}\par
% BEGIN ECONOMICS STATEMENT
\label{problem:OP-0147}
% GAME_THEORY_PROBLEM: {"local_id": "GT-017", "original_number": 17, "kind": "R", "entry_kind": "statement_only", "published": false, "review_required": true, "category": "game_theory"}
\label{game:GT-017}
{\sffamily\small\color{muted}
\textbf{GT-017} | Learning in games and regret minimization\\
\textbf{R} | Editorial minimax specification of the source's fastest-learning agenda.
\par}
\subsection*{Definitions and mathematical statement}
Fix $k\geq2$, action counts $m_1,\ldots,m_k$, and finite games $G$ with payoffs in $[0,1]$. At round $t$, learner $i$ chooses $x_i^t\in\Delta_{m_i}$ using only its own past full-information payoff vectors and its private randomness; no other player's utility function is supplied. Against self-play it receives $v_i^t(a_i)=u_i(a_i,x_{-i}^t)$. Its external regret is
\[
 R_i(T)=\max_{a_i}\sum_{t=1}^Tv_i^t(a_i)-
 \sum_{t=1}^T\langle x_i^t,v_i^t\rangle.
\]
Fix an adversarial-regret budget constant $C_0$ for which admissible learners exist. Let $\mathfrak L$ contain learner tuples whose individual expected positive regret against every nonanticipating sequence $v_i^t\in[0,1]^{m_i}$ is at most $C_0\sqrt{T\log m_i}$ for every $T$, with a declared polynomial per-round computational budget. Define
\[
 \mathcal R_{k,m}(T)=\inf_{\mathcal A\in\mathfrak L}\sup_G
 \max_i\mathbb E[(R_i^{\mathcal A,G}(T))_+],\qquad (r)_+=\max\{r,0\}.
\]
Determine matching asymptotic bounds in $T$, including whether bounded regret is attainable, for this fixed feedback and admissibility model.
\subsection*{Short English statement}
How fast can decentralized learners perform in self-play while keeping guarantees against arbitrary opponents?
\subsection*{Variants and scope boundaries}
Horizon-aware and anytime learners, bandit versus full-information feedback, mixed-payoff versus realized-action regret, and communication permissions change the minimax problem. No-regret self-play controls coarse correlated equilibrium of the empirical joint distribution; it does not generally imply last-iterate Nash convergence.
\subsection*{Source relationship and status}
[S1] asks for the fastest no-regret algorithms in games and proves logarithmic self-play bounds for its framework. The minimax class and the budget $C_0$ above are declared editorial choices rather than a universal minimax formulation attributed to the paper.
\subsection*{References}
\begingroup\footnotesize\raggedright
{\footnotesize\raggedright
\textbf{[S1]} Ashkan Soleymani, Georgios Piliouras, and Gabriele Farina (2025). \emph{Cautious Optimism: A Meta-Algorithm for Near-Constant Regret in General Games}. arXiv:2506.05005v1. Locator: Section 1, fastest-algorithms question; Table 1; Section 2 and Appendix A for regret and feedback. \url{https://arxiv.org/html/2506.05005v1}\par}
\par\endgroup

\subsection*{Common conventions: Source mathematical conventions}

\textbf{Finite games.} For a finite set $A$, $\Delta(A)$ is its probability
simplex. Players choose independent mixed strategies in Nash equilibrium;
correlation is allowed only when explicitly part of the solution concept.
In a game with payoff functions $u_i$ and strategy profile $x$, an additive
$\varepsilon$ Nash equilibrium satisfies
\[
 u_i(x)\geq u_i(x_i',x_{-i})-\varepsilon
 \quad\text{for every player $i$ and every admissible deviation $x_i'$}.
\]
For numerical approximation, payoffs are normalized as locally specified.
An $\varepsilon$ payoff residual is distinct from distance at most
$\varepsilon$ to the set of exact equilibria. Point convergence, convergence
of empirical play, and convergence in distance to an equilibrium set are
also distinct.

\textbf{Computational representations.} Unless stated otherwise, a complexity
question takes rational data with binary encoding; polynomial time is measured
in total bit length. A fixed-accuracy polynomial algorithm is not necessarily
polynomial in $1/\varepsilon$, and neither is necessarily polynomial in
$\log(1/\varepsilon)$. Succinct games and value, demand, or sampling oracles
must declare their interfaces. Exact real solutions require an explicit
representation; an unspecified list of arbitrary real numbers is not a
finite computational output. A claimed algorithmic barrier is meaningful
only for its specified model and reductions.

\textbf{Dynamic games.} Histories, information, observation, and allowable
randomization are part of the model. A stationary strategy depends only on
the current state; a positional strategy in a deterministic graph game
chooses one action at each controlled vertex. Nash equilibrium at an initial
state and equilibrium after every history are different requirements.
An expectation of a pathwise $\liminf$ payoff, a $\liminf$ of expected averages,
a limiting expected average, and a uniform finite-horizon equilibrium are
different criteria. None is silently substituted for another.

\textbf{Allocation and mechanisms.} A complete allocation assigns every item
exactly once unless otherwise stated. Zero-valued goods, negative values,
capacity constraints, feasibility, lotteries, and copies of goods affect
fairness definitions and are specified locally. Dominant-strategy incentive
compatibility, Bayesian incentive compatibility, universal truthfulness,
and truthfulness in expectation impose different inequalities. Whenever
an objective ratio has zero denominator, its convention is stated or those
instances are excluded.

\textbf{Infinite-dimensional models.} Expectations, integrals, stochastic
equations, and optimization objectives require their local regularity,
integrability, filtrations, and admissible domains. A backward--forward
mean-field system is a boundary-value problem; it is not an initial-value
semiflow without an additional construction. Long-time, large-population,
and vanishing-noise limits require an order of limits and a convergence
topology. If a minimum need not be attained, the objective is its infimum
or supremum and the approximation requirement is stated separately.
% END ECONOMICS STATEMENT
\end{document}
